\documentclass[12pt,a4paper]{article}
\usepackage[utf8]{inputenc}
\usepackage[T1]{fontenc}
\usepackage{amsmath,amssymb}
\usepackage{enumitem}
\usepackage{geometry}
\usepackage{array}
\usepackage{multicol}
\geometry{margin=2.2cm}
\setlist{itemsep=0.25em, topsep=0.3em}
\title{Aula 12 -- Processamento de Transações}
\author{Banco de Dados -- Guia de Revisão}
\date{}
\begin{document}
\maketitle
\tableofcontents
\newpage

\section{Objetivo da aula}
Esta aula explica transações, propriedades ACID, estados de transação e planos de execução. Para prova, o ponto mais importante é analisar escalonamentos e decidir se são seriais, serializáveis, restauráveis, sem cascata ou estritos.

\section{Transação}
Uma transação é uma unidade lógica de processamento. Ela deve executar por completo ou não produzir efeito parcial.

Exemplos:
\begin{itemize}
    \item transferência bancária;
    \item compra online;
    \item reserva de passagem;
    \item atualização de estoque.
\end{itemize}

\section{ACID}
\begin{itemize}
    \item \textbf{Atomicidade}: tudo ou nada.
    \item \textbf{Consistência}: leva o banco de um estado consistente a outro.
    \item \textbf{Isolamento}: parece executar isolada das outras.
    \item \textbf{Durabilidade}: depois do \texttt{COMMIT}, o efeito permanece.
\end{itemize}

\section{Comandos SQL}
\begin{verbatim}
START TRANSACTION;
-- comandos
COMMIT;
\end{verbatim}

Se houver erro:
\begin{verbatim}
START TRANSACTION;
-- comandos
ROLLBACK;
\end{verbatim}

\section{Operações de transação}
Modelo simplificado:
\begin{itemize}
    \item $r_i(X)$: transação $T_i$ lê item $X$.
    \item $w_i(X)$: transação $T_i$ escreve item $X$.
    \item $c_i$: transação $T_i$ confirma.
    \item $a_i$: transação $T_i$ aborta.
\end{itemize}

\section{Planos de execução}
Um plano, ou escalonamento, é a ordem em que operações de transações aparecem.

Plano serial:
\[
r_1(X), w_1(X), c_1, r_2(X), w_2(X), c_2
\]
Plano não serial:
\[
r_1(X), r_2(X), w_1(X), w_2(X), c_1, c_2
\]

\section{Serializabilidade}
Um plano é serializável quando produz efeito equivalente a algum plano serial.

\subsection{Conflitos}
Duas operações conflitam quando:
\begin{itemize}
    \item são de transações diferentes;
    \item acessam o mesmo item;
    \item pelo menos uma é escrita.
\end{itemize}

Conflitam:
\begin{itemize}
    \item $r_1(X)$ e $w_2(X)$;
    \item $w_1(X)$ e $r_2(X)$;
    \item $w_1(X)$ e $w_2(X)$.
\end{itemize}

Não conflitam:
\begin{itemize}
    \item $r_1(X)$ e $r_2(X)$;
    \item operações em itens diferentes.
\end{itemize}

\section{Grafo de precedência}
Para verificar serializabilidade por conflito:
\begin{enumerate}
    \item Crie um nó para cada transação.
    \item Para cada conflito, desenhe uma aresta da transação da operação anterior para a posterior.
    \item Se houver ciclo, não é serializável por conflito.
    \item Se não houver ciclo, é serializável por conflito.
\end{enumerate}

\subsection{Exemplo com ciclo}
Plano:
\[
r_1(X),\ w_2(X),\ w_1(X)
\]
Conflitos:
\begin{itemize}
    \item $r_1(X)$ antes de $w_2(X)$ gera $T_1 \rightarrow T_2$.
    \item $w_2(X)$ antes de $w_1(X)$ gera $T_2 \rightarrow T_1$.
\end{itemize}
Há ciclo. Logo, não é serializável por conflito.

\section{Restaurabilidade}
Um plano é restaurável se uma transação só confirma depois das transações das quais leu dados confirmarem.

Não restaurável:
\[
w_1(X),\ r_2(X),\ c_2,\ a_1
\]
$T_2$ leu valor de $T_1$ e confirmou antes de saber se $T_1$ confirmaria.

\section{Plano sem cascata}
Um plano sem cascata só permite leitura de valores já confirmados.

Não sem cascata:
\[
w_1(X),\ r_2(X),\ c_1,\ c_2
\]
Mesmo sendo restaurável, $T_2$ leu $X$ antes do \texttt{COMMIT} de $T_1$.

Sem cascata:
\[
w_1(X),\ c_1,\ r_2(X),\ c_2
\]

\section{Plano estrito}
Em plano estrito, se uma transação escreve $X$, nenhuma outra pode ler nem escrever $X$ até a primeira confirmar ou abortar.

Violação:
\[
w_1(X),\ w_2(X),\ c_1,\ c_2
\]
$T_2$ escreve $X$ antes de $T_1$ confirmar.

\section{Problemas clássicos}
\begin{itemize}
    \item \textbf{Leitura suja}: ler valor escrito por transação ainda não confirmada.
    \item \textbf{Atualização perdida}: duas transações atualizam o mesmo valor e uma sobrescreve o efeito da outra.
    \item \textbf{Plano não restaurável}: uma transação confirma baseada em valor de outra que depois aborta.
\end{itemize}

\section{Pegadinhas comuns}
\begin{itemize}
    \item Serial e serializável não são a mesma coisa.
    \item Todo plano serial é serializável, mas nem todo serializável é serial.
    \item Plano sem cascata é mais forte que restaurável.
    \item Plano estrito é ainda mais restritivo.
\end{itemize}

\section{Desafios}
\subsection*{Desafio 1}
Analise:
\[
r_1(X),\ w_2(X),\ w_1(X)
\]
O grafo de precedência tem ciclo?

\subsection*{Desafio 2}
Analise:
\[
w_1(X),\ r_2(X),\ c_2,\ a_1
\]
O plano é restaurável?

\section{Resolução dos desafios}
\textbf{Desafio 1}: sim. Há $T_1\rightarrow T_2$ e $T_2\rightarrow T_1$. Não é serializável por conflito.

\textbf{Desafio 2}: não. $T_2$ leu valor escrito por $T_1$ e confirmou antes de $T_1$. Se $T_1$ aborta, $T_2$ confirmou com dado inválido.

\section{Checklist final}
\begin{itemize}
    \item Sei definir ACID?
    \item Sei reconhecer conflito entre operações?
    \item Sei montar grafo de precedência?
    \item Sei diferenciar restaurável, sem cascata e estrito?
    \item Sei identificar leitura suja e atualização perdida?
\end{itemize}

\end{document}
